\color{unidarkgreen}{\subsection{Контроль цепей напряжения}}
\color{black}

\begin{enumerate}[label=\arabic{section}.\arabic{subsection}.\arabic*,  align=left, labelsep=4pt, leftmargin=0pt, itemindent=57pt]

\item
Контроль цепей напряжения предназначен для выявления нарушений в цепях напряжения низкой
и средней сторон трансформатора. На рисунке \ref{pic:KCN} представлен контроль цепей напряжения НН (для НН1 и
НН2 аналогично).

\begin{figure}[H]
\centering
\includegraphics[width=0.9\textwidth,height=0.9\textheight,keepaspectratio]{\fbpath/081.1741 КЦН/img/KCN.pdf}
\captionof{figure}{Функциональный блок КЦН}\label{pic:KCN}
\end{figure}

\item
Цепи напряжения подключаются к ТН, установленным на шинах НН и СН.

\item
Контроль цепей напряжения низкой и средней стороны трансформатора основан на выявлении
нарушения симметрии в цепях напряжения и просадки линейного напряжения соответствующей
стороны.

\item
Ввод в работу КЦН осуществляется независимо для каждой стороны с помощью уставок \textbf{«КЦН
НН», «КЦН НН1» , «КЦН НН2»}. Оперативный ввод/вывод КЦН производится с помощью БУ (блоков
управления функциями) «Вывод КЦН НН», «Вывод КЦН НН1», «Вывод КЦН НН2».

\item
Порог срабатывания минимального линейного напряжения задается уставками \textbf{Uл КЦН НН»}, \textbf{Uл
КЦН НН1»}, \textbf{Uл КЦН НН2»}. Порог срабатывания максимального напряжения обратной
последовательности задается уставками  \textbf{«U2 КЦН НН»}, \textbf{«U2 КЦН НН1»}, \textbf{«U2 КЦН НН2»}. Выдержка
времени контроля неисправности цепей напряжения задается уставками \textbf{«Тср КЦН НН»}, \textbf{«Тср КЦН
НН1»}, \textbf{«Тср КЦН НН2»}.

\item
Уставки функции КЦН представлены в таблице \ref{desc:tbl1}.

\small
\begin{longtable}{|>{\centering\arraybackslash}m{5.3cm}|>{\centering\arraybackslash}m{4.2cm}|>{\centering\arraybackslash}m{1cm}|>{\centering\arraybackslash}m{3.3cm}|>{\centering\arraybackslash}m{1.8cm}|}
% Первая страница: полный заголовок
\caption{Перечень уставок функции КЦН\hfill\vspace{-0.5\baselineskip}}\label{desc:tbl1}\\
\hline
\rowcolor{gray!30}
Наименование уставки & Диапазон & Шаг & Значение по умолчанию & Единица измерения \\
\hline
\endfirsthead

% Промежуточные страницы: «Продолжение таблицы»
\caption*{\hspace{3pt}\emph{Продолжение таблицы \ref{desc:tbl1}}\hfill\ }\vspace{-0.5\baselineskip} \\
\hline
\rowcolor{gray!30}
Наименование уставки & Диапазон & Шаг & Значение по умолчанию & Единица измерения \\
\hline
\endhead

%===m>ДЗАТ|-
%===m
\end{longtable}
\normalsize	

\end{enumerate}